\section{État de l'art}
\label{sec:etat}

\textbf{Quatre courants de recherche} touchent à la préconisation de décisions face aux perturbations : la mesure de la résilience, les stratégies d'atténuation et de contingence, la simulation des perturbations et les réseaux bayésiens causaux. Leur revue détaillée figure en section~S2 ; on en retient ici ce qui situe la démarche et permet l'appropriation des apports de ce travail.

Le premier courant mesure la résilience à partir de la trajectoire de performance. La résilience d'une chaîne logistique est généralement définie comme sa capacité à se préparer aux perturbations, à y répondre et à s'en rétablir \citep{tukamuhabwa2015, kamalahmadi2016}. \citet{rahiel2026in4pl} et \citet{rahiel2026these} ont formalisé mathématiquement cette définition à partir de la trajectoire de performance, qui décrit la chute puis le rétablissement d'une chaîne perturbée. Ce cadre a été appliqué à la chaîne logistique hospitalière \citep{rahiel2025esd} et complété par une modélisation probabiliste des perturbations sévères \citep{rahiel2026codit}. Plusieurs mesures partent ainsi de la performance observée dans le temps. \citet{bruneau2003} associent la perte de résilience à la surface comprise entre la performance nominale et la performance dégradée, \citet{henry2012} traitent la résilience comme une fonction du temps et d'autres en tirent des indicateurs utiles à la décision \citep{spiegler2012, behzadi2020}. Tout en intégrant tout cela, la formalisation de \citet{rahiel2026in4pl} décrit la courbe entière par une fonction composée de tangentes hyperboliques à sept paramètres et qui fournit six indicateurs. Il faut noter toute fois que l'approche suppose une courbe déjà observée. Aucune ne dit comment obtenir en nombre des courbes comparables, pour des crises variées et des réponses différentes. 

Mesurer la résilience ne dit pas comment l'améliorer ; c'est l'objet des travaux sur les stratégies d'atténuation et de contingence. On distingue les mesures d'atténuation, prises avant la perturbation, des mesures de contingence, prises une fois qu'elle survient \citep{tang2006, tomlin2006}. Les travaux quantitatifs ont précisé le rôle de chaque levier, le stock agissant plutôt en prévention et la capacité de réserve en réaction \citep{lucker2019}, comparé politiques proactives et politiques de rétablissement \citep{ivanov2016tre} et chiffré les coûts de la résilience \citep{aldrighetti2021}. Dans ces travaux, les décisions sont le plus souvent fixées à l'avance. Leur déclenchement au fil de la crise, avec les délais de détection et de décision d'un responsable des opérations réel, est peu modélisé \citep{klibi2012ijpe}. On compare rarement plusieurs politiques de gestion face à un même aléa.

Pour confronter ces stratégies à des perturbations variées, la simulation s'est imposée. Elle est devenue l'outil courant pour étudier la propagation des perturbations, ou effet domino notamment dans la supply chain \citep{schmitt2012, ivanov2017}. Elle sert aussi à construire l'approche par plans d'expériences \citep{macdonald2018}, à générer des scénarios plausibles \citep{klibi2012ejor} et à comparer des stratégies de rétablissement \citep{ramzy2022}. ISOMORPH, premier jumeau numérique public d'un réseau logistique multi-échelon à notre connaissance, fournit des trajectoires de référence pour la prévision, sans décision qui réagisse à l'état du réseau \citep{zhang2026isomorph}. Ce simulateur est conçu pour étudier un réseau ou un cas. Aucun simulateur ne rejoue systématiquement un même aléa sans réaction puis sous plusieurs politiques de gestion.

Reste à tirer des crises observées ou simulées des règles de décision, ce que permettent les réseaux bayésiens. Ils représentent les dépendances entre risques, stratégies et conséquences et l'on peut y raisonner de la cause vers l'effet comme de l'effet vers la cause \citep{hosseiniivanov2020}. Ils ont servi à modéliser la propagation des risques \citep{garvey2015, ojha2018} et à choisir des stratégies d'atténuation \citep{qazi2018}. Leur structure, souvent fixée par des experts, peut aussi être apprise sur des données et les deux sources se combinent \citep{heckerman1995}. Ce dernier point représente un avantage remarquable. Pour qu'un tel modèle préconise une décision, l'effet de cette décision doit être identifiable. Une décision est une intervention au sens de \citet{pearl2009} et sa conséquence ne se lit pas dans une corrélation observée, puisque le responsable des opérations réagit d'autant plus fort que la crise est grave. Les historiques d'entreprise ne permettent pas cette identification de nature causale, car souvent ils ne contiennent ni l'étiquette de la décision prise ni ce qui se serait passé sans elle \citep{do2024}.

Ces quatre courants sont solides, mais ils se sont développés séparément (tableau~\ref{tab:synthese}). Aucune base publique connue ne relie, pour un grand nombre de crises, la perturbation subie, la décision prise et la trajectoire du taux de service qui en découle.

\begin{table*}[htbp]
\caption{Courants de la littérature, limites pour la préconisation de décisions et réponse apportée par ce travail.}
\label{tab:synthese}
\footnotesize
\begin{tabularx}{\textwidth}{@{}P{1.9cm}P{3.3cm}LLL@{}}
\toprule
\textbf{Courant} & \textbf{Travaux représentatifs} & \textbf{Apport} & \textbf{Limite pour la préconisation} & \textbf{Réponse de ce travail} \\
\midrule
Mesure de la résilience
  & \citealp{bruneau2003}; \citealp{henry2012}; \citealp{spiegler2012}; \citealp{behzadi2020}; \citealp{rahiel2026in4pl}
  & Courbe de performance, indicateurs, fonction explicite de la trajectoire
  & Suppose une courbe déjà observée
  & Mesurer : fonction à sept paramètres ajustée sur chaque trajectoire simulée \\
\addlinespace
Atténuation et contingence
  & \citealp{tang2006}; \citealp{tomlin2006}; \citealp{lucker2019}; \citealp{ivanov2016tre}; \citealp{aldrighetti2021}
  & Répertoire de leviers et de leurs coûts
  & Décisions fixées à l'avance ; délais de réaction peu modélisés
  & Simuler : responsable des opérations qui ne voit que le taux de service, trois profils de gestion \\
\addlinespace
Simulation des perturbations
  & \citealp{schmitt2012}; \citealp{ivanov2017}; \citealp{macdonald2018}; \citealp{klibi2012ejor}; \citealp{ramzy2022}; \citealp{zhang2026isomorph}
  & Propagation des ruptures, génération de scénarios, stratégies de rétablissement
  & Étude d'un cas ; pas de contrefactuels systématiques
  & Simuler : chaque crise rejouée sans réaction puis sous chaque profil \\
\addlinespace
Réseaux bayésiens causaux
  & \citealp{garvey2015}; \citealp{ojha2018}; \citealp{qazi2018}; \citealp{hosseiniivanov2020}; \citealp{rahiel2026codit}
  & Inférence dans les deux sens, choix de stratégies
  & Données rares, sans étiquette de décision ni contrefactuel
  & Apprendre : réseau appris sur la base contrefactuelle, avec arcs interdits \\
\bottomrule
\end{tabularx}
\end{table*}
